\section{Reconstructing the exposure}
\label{sec:reconstruction}

This section sets out the construction. It is specification rather than result, and
its purpose is to make the exercise repeatable and to put the judgement calls where a
reader can see them. The diagnostics quoted here are reported in full in
Appendices~\ref{app:curve} and~\ref{app:pricing}.

\subsection{From filings to a bucketed balance sheet}

Both banks report securities by type and by contractual maturity band, deposits split
between insured and uninsured, loans by product, and borrowings by term. Those tables
are parsed from the inline XBRL of each filing into 192 positions across the four
bank-years, each carrying an amortised cost, a fair value where one is disclosed, a
contractual maturity bucket and a book yield.

Two parsing conventions matter enough to state. An em-dash in these tables means a
disclosed zero, while an empty cell is layout padding; conflating the two silently
shifts a row's figures into the wrong columns, which is a failure mode that produces
plausible numbers rather than an error. And the maturity tables are not reported on a
consistent basis between the two banks: SVB's held-to-maturity bands are given at
net carrying amount while Schwab's are given at fair value. Each bank's basis is
therefore read from the table it belongs to rather than assumed. Totals are
reconciled against the balance sheet figures in the same filing before anything is
valued.

Positions are priced at a representative maturity for their bucket: six months for
under one year, three years for one to five, seven and a half for five to ten, and
twenty years for over ten. The open-ended final bucket is the consequential one,
since it holds the longest and most rate-sensitive paper and the filings do not say
how far beyond ten years it reaches. Section~\ref{sec:sensitivity} varies this
assumption rather than defending it.

\subsection{The curve}

The U.S. Treasury publishes par yields, not spot rates \citep{Treasury}. Discounting
cash flows directly with par yields is a standard way for a reconstruction like this
one to go
quietly wrong, because the error grows with maturity and with the slope of the curve
--- which is exactly where SVB's book sat.

Zero rates are therefore bootstrapped from the published par curve, solving each node
so that a par bond of that maturity prices to one hundred, with intermediate coupon
dates interpolated against the trial rate so the construction is self-consistent. The
resulting zeros are then fitted with the Svensson extension of the Nelson--Siegel
form \citep{NelsonSiegel1987, Svensson1994}, which gives a smooth curve with an
analytic forward rate --- needed because the short-rate model is fitted to the forward
curve rather than to the zeros.

The fit is scored against par-yield error rather than against the bootstrapped zeros
it was fitted through, which is the comparison a reviewer would make and the harder of
the two. The worst residual is 8.70 basis points at the end of 2021 and 10.05 at the
end of 2022, both at the twenty-year point, with root-mean-square errors of 3.88 and
5.39 basis points. Fitting is multi-start with economic bounds on the level
parameter, because the unconstrained problem admits degenerate solutions that fit the
observed points acceptably while implying negative forward rates at long horizons.

\subsection{The short rate}

Securities with prepayment options cannot be valued off a curve alone, so the short
rate follows a one-factor Hull--White process \citep{HullWhite1990},

\begin{equation}
\mathrm{d}r(t) = \bigl(\theta(t) - a\,r(t)\bigr)\,\mathrm{d}t + \sigma\,\mathrm{d}W(t),
\label{eq:hw}
\end{equation}

which is the mean-reverting specification of \citet{Vasicek1977} with a
time-dependent drift chosen so that the model reproduces the initial term structure
exactly. Internally the short rate is decomposed as $r(t) = x(t) + \alpha(t)$, with
$x$ a driftless mean-reverting Gaussian starting at zero and $\alpha$ a deterministic
function carrying the initial curve. That decomposition lets the simulation be
stepped exactly, with no discretisation bias in the short rate itself, and keeps the
fit to the initial curve exact however coarse the time grid.

Mean reversion and volatility are estimated from the realised path of the three-month
bill \citep{FRED}, because a swaption surface is not public. Monthly observations from
1990 give a
mean-reversion half-life of 7.80 years and a volatility of 61.87 basis points. That
specification was chosen rather than assumed: the same estimator run daily over
2003--2022 returns a half-life of 365 years and a long-run mean of 86\%, which is a
regression failing to identify anything rather than a fact about interest rates.
Section~\ref{sec:sensitivity} reports the full grid and is candid that mean reversion
is weakly identified under every specification tried.

Two consequences follow from calibrating historically rather than to traded options.
The parameters carry no market price of interest-rate risk, which will tend to
undervalue the prepayment option embedded in a mortgage book; and the volatility is
an estimate of realised variation rather than of implied. Neither is corrected for.
The option-adjusted spread absorbs what is left, and
Section~\ref{sec:sensitivity} establishes that the rate parameters are not where the
residual error lives.

Valuation is by Monte Carlo with antithetic draws, 20,000 paths and monthly steps,
and the same draws are reused across scenarios so that differences between scenarios
measure curvature rather than sampling noise. Simulated discount factors reproduce
the closed-form Hull--White bond price to within 0.12 basis points of yield at every
tested maturity, which is the gate on everything downstream.

\subsection{Prepayment and option-adjusted spreads}

Agency pass-throughs prepay, and prepayment is modelled as an S-curve in refinancing
incentive --- the gap between the pool's weighted-average coupon and the prevailing
mortgage rate --- running from a floor of 5\% to a ceiling of 45\% conditional
prepayment rate, centred at the money and turning over roughly two percentage points
of incentive either side. The mortgage rate is the model's own ten-year zero rate plus
a spread anchored at each valuation date to the published Freddie Mac survey rate
\citep{FreddieMac}.

These parameters are stylised. They are not estimated from pool-level data, because
the filings disclose neither loan-level balances nor actual weighted-average coupons,
and they are held fixed across both banks and both valuation dates so that nothing can
be tuned after the answer is visible.

What is calibrated is one option-adjusted spread per security-type and maturity
bucket, solved so that the model reproduces each bank's disclosed fair value at
31~December 2021. Thirteen spreads for SVB and ten for Schwab, ranging from $-99$ to
$87$ basis points across SVB's fitted buckets and $-168$ to $110$ across Schwab's.
Twelve of SVB's thirteen and all ten of Schwab's solve to the disclosed value
essentially exactly.

The exception should be named. SVB's \$100m holding of variable-rate agency CMOs is
disclosed above par, at \$101m, and a model that discounts at a spread cannot reach
that level for a floating-rate position; the solver pins the spread at its $-300$
basis point lower bound and still returns \$100.73m. The position is 0.08\% of SVB's
securities book and the shortfall is \$0.27m, so it does not affect any figure
reported in this paper, but it is a bucket the calibration failed to fit rather than
one it fitted well.

The spreads are then frozen. Nothing about them is re-estimated at the 2022 valuation
date, which is what makes the test in Section~\ref{sec:comparison} out of sample.

\subsection{The deposit franchise}

Non-maturity deposits are valued as a below-market funding source that decays, in the
sense of \citet{Drechsler2021}: a deposit base paying less than market rates is worth
less than its face amount to the depositor and correspondingly more to the bank, and
that value is an asset which rising rates increase and departing depositors destroy.
Writing $D_0$ for the book balance, $h$ for the half-life of the balance, $\beta$ for
the pass-through from market to deposit rates, $y_t$ for the market rate and $c$ for a
servicing cost,

\begin{equation}
V = \sum_t D_0 \, e^{-t \ln 2 / h} \,
    \bigl[(1 - \beta)\, y_t - c\bigr] \, \mathrm{DF}_t ,
\label{eq:deposits}
\end{equation}

with franchise value taken as the book balance less $V$. Deposits are split into an
insured and an uninsured layer, each with its own half-life, because the uninsured
layer is the one that can leave.

Three parameters, of which one is disclosed. SVB's FY2021 filing states a deposit
beta of 0.60, so that is used rather than estimated \citep{SVB10K2021}. The servicing
cost is set at 35 basis points, within the usual range and not a quantity the result
turns on. The half-lives are disclosed by nobody, and they are where the exercise
either stands or falls; Sections~\ref{sec:comparison} and~\ref{sec:sensitivity} treat
them accordingly.

\subsection{Scenarios}

The shocks are the six standardised scenarios of the Basel Committee's 2016 standard
in its USD calibration \citep{BCBS2016}: 200 basis points parallel, 300 short and 150
long, with the short and long shapes decaying over four years, and the two rotation
scenarios built from weighted combinations of them. The supervisory outlier threshold
is 15\% of Tier~1 capital.

Two implementation points change the answer. The shock applies to the zero curve, so
the instantaneous forward curve has to be re-derived rather than shifted by the same
amount; only the parallel case leaves the two equal, and the short-rate model is
fitted to forwards. And post-shock rates are floored at zero, as the standard
specifies. At the end of 2021 the entire front of the curve sat below 2\%, so a
downward 200 basis point shock is substantially absorbed by the floor and the realised
down-shock is nothing like its nominal size. A model that omits the floor reports a
near-symmetric profile for a balance sheet whose exposure was anything but, and would
miss the non-monotonicity that Section~\ref{sec:comparison} treats as its main
qualitative result.
